\uuid{uJXC}
\titre{Limite}
\auteur{bodin ; hanani ; mzari}
\organisation{exo7}
\licence{CC BY-NC-SA 4.0}
\datecreate{2020-01-20}
\niveau{L1}
\module{Analyse}
\chapitre{Suite}
\sousChapitre{Autre}
\difficulte{2}
\theme{Suites, limite}
\exo7qcmid{lille-1.suites.014}

\contenu{
\begin{question}
\qclassification{121.00}
\qtags{motcle=limite}
Soient $\displaystyle u_n=\sqrt{n^2+4n-1}-n$ et $\displaystyle v_n=\frac{4n-1}{\sqrt{n^2+4n-1}+n}$. Quelles sont les bonnes réponses ?
\begin{answers}
    \bad{$\displaystyle \lim _{n\to +\infty}u_n=+\infty$ et $\displaystyle \lim _{n\to +\infty}v_n=0$ }
    \bad{$\displaystyle \lim _{n\to +\infty}u_n=0$ et $\displaystyle \lim _{n\to +\infty}v_n=2$}
    \good{$\displaystyle \lim _{n\to +\infty}u_n=2$ et $\displaystyle \lim _{n\to +\infty}v_n=2$}
    \bad{$\displaystyle \lim _{n\to +\infty}u_n=+\infty$ et $\displaystyle \lim _{n\to +\infty}v_n=0$}
\end{answers}
\begin{explanations}
On multiplie par le terme conjugué, on obtient $\displaystyle u_n=\frac{4n-1}{\sqrt{n^2+4n-1}+n}=v_n$. Ensuite,
$$\frac{4n-1}{\sqrt{n^2+4n-1}+n}=\frac{4n\left(1-\frac{1}{4n}\right)}{n\left(\sqrt{1+\frac{4}{n}-\frac{1}{n^2}}+1\right)}=4\frac{1-\frac{1}{4n}}{\sqrt{1+\frac{4}{n}-\frac{1}{n^2}}+1}\underset{+\infty}{\longrightarrow}2.$$
\end{explanations}
\end{question}
}
